\chapter{Expérimentations et résultats}
\label{chap:experimentation}

\section*{Introduction}
\addcontentsline{toc}{section}{Introduction}
Ce chapitre évalue HAKE-MER sur GoEmotions par une ablation
cumulative à protocole strictement constant, complétée par une
comparaison des ressources de connaissances externes et par des
analyses qualitative et d'erreur. Les hypothèses H1 à H4 du
chapitre 3 y sont confrontées aux résultats.

\section{Cadre expérimental}
\label{sec:expe-cadre}

\subsection{Jeu de données}
Toutes les expériences utilisent GoEmotions \cite{demszky2020}
(configuration \emph{simplified}) : 43\,410 exemples d'entraînement,
5\,426 de validation et 5\,427 de test, 27 émotions plus la classe
neutre ($K{=}28$), en multi-étiquettes. L'analyse exploratoire
éclaire trois choix : le \textbf{déséquilibre} (\emph{neutral} :
14\,219 occurrences contre 77 pour \emph{grief}, rapport de 185)
impose le F1-macro comme métrique décisive ; la
\textbf{multi-étiquette} est réelle (16{,}4\,\% des exemples portent
au moins deux émotions ; paires fréquentes :
\emph{admiration}+\emph{gratitude} 279,
\emph{anger}+\emph{annoyance} 269) ; la \textbf{structure} (longueur
moyenne 19 tokens, 42{,}2\,\% de textes multi-phrases) justifie les
bornes du prétraitement hiérarchique ($M{=}4$ phrases de 32 tokens —
à 20 tokens, 24{,}1\,\% des exemples étaient tronqués contre
1{,}8\,\% à 32).

\subsection{Protocole}
Toutes les configurations partagent : 5~époques (budget fixé après
constat d'un plateau au-delà), graine aléatoire 42, taille de lot 24
(32 pour la ligne de base plate), précision mixte sur GPU~T4, et les
mêmes métriques : F1-micro, \textbf{F1-macro} (métrique de référence,
seule sensible aux émotions rares), correspondance exacte et mAP.
L'architecture est implémentée comme une classe unique à indicateurs
d'activation : les configurations comparées ne diffèrent que par les
modules activés.

\section{Choix des backbones}
\label{sec:expe-backbones}

Un comparatif préliminaire de quatre encodeurs à protocole identique
(tableau~\ref{tab:backbones}) a guidé le choix : RoBERTa-base
\cite{liu2019roberta} domine sur F1-macro et mAP ; DistilBERT-base
\cite{sanh2019distilbert} offre le meilleur compromis
coût/performance ; ALBERT \cite{lan2020albert}, malgré son faible
nombre de paramètres, est le plus lent à entraîner et le plus faible
en F1-macro. La suite rapporte systématiquement DistilBERT \emph{et}
RoBERTa.

\begin{table}[H]
\centering
\small
\begin{tabular}{lccccc}
\toprule
Modèle & Params & Temps (5 ép., T4) & F1-micro & F1-macro & mAP\\
\midrule
RoBERTa-base \cite{liu2019roberta}       & 125M & 805\,s  & \textbf{0,588} & \textbf{0,488} & \textbf{0,516}\\
BERT-base \cite{devlin2019}              & 110M & 787\,s  & 0,579 & 0,480 & 0,500\\
DistilBERT-base \cite{sanh2019distilbert}& 67M  & \textbf{415\,s} & 0,579 & 0,475 & 0,498\\
ALBERT-base-v2 \cite{lan2020albert}      & 12M  & 1\,020\,s & 0,584 & 0,467 & 0,491\\
\bottomrule
\end{tabular}
\caption{Comparatif préliminaire des encodeurs sur GoEmotions (test),
protocole identique.}
\label{tab:backbones}
\end{table}

\section{Ablation cumulative}
\label{sec:expe-ablation}

Le tableau~\ref{tab:ablation} présente le résultat central du
travail : l'ablation cumulative des quatre modules, sur les deux
backbones.

\begin{table}[H]
\centering
\small
\begin{tabular}{l cccc cccc}
\toprule
& \multicolumn{4}{c}{RoBERTa-base} & \multicolumn{4}{c}{DistilBERT-base}\\
\cmidrule(lr){2-5}\cmidrule(lr){6-9}
Configuration & F1-mi & F1-ma & Exact & mAP & F1-mi & F1-ma & Exact & mAP\\
\midrule
Ligne de base (PLM affiné)      & 0,584 & 0,476 & \textbf{0,475} & 0,504 & \textbf{0,590} & 0,459 & \textbf{0,473} & 0,502\\
M1 (hiérarchique)               & 0,575 & 0,509 & 0,461 & 0,512 & 0,558 & 0,438 & 0,448 & 0,487\\
M1+M2 (attention croisée)       & \textbf{0,588} & 0,510 & 0,466 & 0,520 & 0,586 & 0,480 & 0,468 & 0,505\\
M1+M2+M3 (enrichissement)       & 0,585 & \textbf{0,523} & 0,458 & \textbf{0,532} & 0,586 & \textbf{0,499} & 0,467 & 0,511\\
M1+M2+M3+M4 (complet)           & 0,586 & 0,496 & 0,461 & 0,520 & 0,584 & 0,485 & 0,463 & \textbf{0,512}\\
\bottomrule
\end{tabular}
\caption{Ablation cumulative de HAKE-MER sur GoEmotions (test),
protocole constant (5~époques, graine~42). Les chiffres RoBERTa
proviennent d'une unique session entraînant les cinq configurations
(modèle final) ; les chiffres DistilBERT de runs par module au même
budget (meilleur point de contrôle sur validation).}
\label{tab:ablation}
\end{table}

\subsection{Lecture des résultats}

\textbf{Le F1-macro progresse à chaque module jusqu'à M3.} Sur
RoBERTa : $0{,}476 \to 0{,}509$ (M1) $\to 0{,}510$ (M2)
$\to \mathbf{0{,}523}$ (M3), soit $\mathbf{+0{,}047}$ au-dessus de la
ligne de base — sur la métrique qui capture les émotions rares,
l'objectif central du sujet. Le mAP suit la même trajectoire (0,504
$\to$ 0,532). L'hypothèse H1 est validée \emph{conditionnellement au
backbone} : la hiérarchie seule améliore nettement RoBERTa
($+0{,}033$ de F1-macro) mais \emph{dégrade} DistilBERT ($-0{,}021$).

\textbf{L'attention croisée corrige la hiérarchie seule (H2).} Sur
les deux backbones, M2 restaure le F1-micro que M1 dégradait
(DistilBERT : $0{,}558 \to 0{,}586$) tout en maintenant ou améliorant
le F1-macro — cohérent avec les co-occurrences mesurées.

\textbf{Le module 4 ne justifie pas son ajout (H4).} La pondération
dynamique fait reculer le F1-macro sur les deux backbones
($0{,}523 \to 0{,}496$ sur RoBERTa). Ce résultat négatif, assumé,
illustre la nécessité de l'ablation contrôlée : une composante
plausible peut ne rien apporter. La meilleure configuration retenue
est donc \textbf{M1+M2+M3}.

\section{Comparaison des connaissances externes}
\label{sec:expe-ressources}

L'hypothèse H3 est testée à architecture et budget fixés, en ne
faisant varier que la ressource injectée par la porte du module~3
(tableau~\ref{tab:ressources}).

\begin{table}[H]
\centering
\small
\begin{tabular}{llccccc}
\toprule
Backbone & Ressource (type) & F1-micro & F1-macro & Exact & mAP & porte $g$\\
\midrule
DistilBERT & NRC (lexical, 10-d)          & 0,586 & \textbf{0,499} & 0,467 & \textbf{0,511} & $+0{,}029$\\
DistilBERT & SenticNet (conceptuel, 5-d)  & 0,587 & 0,486 & 0,470 & 0,507 & $+0{,}027$\\
RoBERTa    & NRC (lexical, 10-d)          & \textbf{0,595} & \textbf{0,492} & \textbf{0,479} & \textbf{0,525} & $-0{,}027$\\
RoBERTa    & SenticNet (conceptuel, 5-d)  & 0,588 & 0,461 & 0,471 & 0,510 & $+0{,}023$\\
\bottomrule
\end{tabular}
\caption{Module 3 avec deux ressources de connaissance, à
architecture et budget fixés. La porte $g$ non nulle atteste
l'utilisation effective de la connaissance.}
\label{tab:ressources}
\end{table}

Le NRC \cite{mohammad2013} domine SenticNet
\cite{cambria2018senticnet} sur le F1-macro pour les deux backbones
($+0{,}013$ DistilBERT, $+0{,}031$ RoBERTa) : les catégories
d'émotions du lexique s'alignent directement sur les étiquettes de
GoEmotions, là où les dimensions abstraites de SenticNet exigeraient
une mise en correspondance que l'injection par biais ne fournit pas.
H3 est validée : le rendement de la connaissance dépend de
l'\emph{alignement taxonomique}.

\section{Analyse qualitative}
\label{sec:expe-qualitatif}

Le suivi d'exemples de classes rares, configuration par
configuration (tableau~\ref{tab:qualitatif}), illustre le rôle de
chaque module : la hiérarchie corrige \emph{relief} là où la ligne de
base prédit \emph{sadness} ; l'attention croisée retrouve le couple
complet \emph{joy}+\emph{surprise}. À l'inverse, \emph{grief} n'est
détectée par aucune configuration (6 exemples de test seulement) —
l'illustration extrême du problème des classes quasi vides.

\begin{table}[H]
\centering
\footnotesize
\begin{tabular}{p{4.6cm}p{2.2cm}p{2.0cm}p{2.0cm}p{2.4cm}}
\toprule
Exemple (extrait) & Vraies étiquettes & Ligne de base & M1 & M1+M2(+M3)\\
\midrule
« Resetting a dislocated knee hurts... feels a lot better immediately after » & relief & sadness & \textbf{relief} & \textbf{relief}\\
« A surprise turn of events! I'm so glad... » & joy, surprise & surprise & admiration & \textbf{joy, surprise}\\
« [NAME] death is just so... senseless. Why? WHY??? » & grief & sadness & disappointment & sadness\\
\bottomrule
\end{tabular}
\caption{Évolution des prédictions sur des exemples de classes rares
(backbone RoBERTa).}
\label{tab:qualitatif}
\end{table}

\section{Analyse d'erreur}
\label{sec:expe-erreur}

Sur le modèle complet (RoBERTa), les confusions sont dominées par la
frontière avec \emph{neutral} : les paires les plus fréquentes
(étiquette vraie manquée $\to$ étiquette prédite à la place) sont
neutral$\to$curiosity (69), annoyance$\to$neutral (64) et
approval$\to$neutral (63) — le modèle peine à départager les émotions
\emph{subtiles} du neutre, davantage qu'à confondre deux émotions
marquées. Le F1 par classe confirme la structure attendue : minimal
sur \emph{grief} (0,00), \emph{relief} (0,13), \emph{realization}
(0,26) et \emph{nervousness} (0,30), maximal sur \emph{gratitude}
(0,92), \emph{amusement} (0,82) et \emph{love} (0,80) — classes
fréquentes aux indices lexicaux saillants.

\section{Positionnement par rapport à la littérature}
\label{sec:expe-sota}

À titre indicatif, le F1-macro obtenu par la meilleure configuration
(M1+M2+M3, RoBERTa) est rapproché du chiffre de référence publié par
les auteurs du jeu de données GoEmotions \cite{demszky2020}, qui
rapportent un F1-macro de 0,46 et un F1-micro de 0,51 pour un
classifieur BERT-base affiné sur la taxonomie complète à 28 classes
(tableau~\ref{tab:sota}).

\begin{table}[H]
\centering
\small
\begin{tabular}{llcc}
\toprule
Étude & Modèle & F1-micro & F1-macro\\
\midrule
Demszky et al. (2020) \cite{demszky2020} & BERT-base & 0,51 & 0,46\\
Notre travail (ce mémoire) & HAKE-MER M1+M2+M3 (RoBERTa) & \textbf{0,585} & \textbf{0,523}\\
\bottomrule
\end{tabular}
\caption{Repère indicatif par rapport au chiffre de référence publié
sur GoEmotions. Comparaison non strictement contrôlée : le protocole
d'entraînement original (nombre d'époques, taille de séquence,
recherche d'hyperparamètres) n'est pas reproduit à l'identique ; seule
la comparaison interne par ablation (tableau~\ref{tab:ablation}) isole
rigoureusement l'apport des modules de HAKE-MER.}
\label{tab:sota}
\end{table}

Cet écart positif est cohérent avec les gains internes mesurés par
ablation, mais ne constitue pas une preuve de supériorité au sens
strict : contrairement au protocole d'ablation cumulative de ce
chapitre, où seul le budget d'entraînement est maintenu constant, la
configuration exacte de \cite{demszky2020} (nombre d'époques, longueur
de séquence, recherche d'hyperparamètres) n'a pas été reproduite à
l'identique. Ce repère situe néanmoins HAKE-MER par rapport au point
d'ancrage le plus largement cité de la littérature sur GoEmotions,
sans remplacer la comparaison contrôlée qui reste la preuve principale
de ce mémoire.

\section{Leçons méthodologiques}
\label{sec:expe-lecons}

Trois incidents résolus en cours de développement méritent d'être
consignés, car chacun modifiait matériellement les conclusions :

\begin{itemize}
\item \textbf{Troncature du prétraitement.} La première borne de
  20 tokens/phrase tronquait 18{,}7\,\% des exemples et faisait
  passer le module 1 \emph{sous} la ligne de base — l'écart mesurait
  le prétraitement, pas l'architecture. Borne recalibrée à 32 tokens.
\item \textbf{Échelle d'initialisation.} Les requêtes d'émotions
  initialisées en $\mathcal{N}(0,1)$ produisaient une perte initiale
  de 1,86 (contre 0,18 après correction en
  $\mathcal{N}(0,0{,}02^2)$).
\item \textbf{Équité des budgets.} Comparer des configurations à
  budgets d'entraînement différents inversait certaines conclusions ;
  d'où le protocole à budget strictement constant.
\end{itemize}

\section*{Conclusion}
\addcontentsline{toc}{section}{Conclusion}
L'ablation cumulative établit que la configuration M1+M2+M3 sur
RoBERTa atteint un F1-macro de 0,523, soit $+0{,}047$ au-dessus du
PLM affiné, le gain se concentrant sur les émotions rares ; le NRC
surpasse SenticNet à architecture fixée ; et le module 4 ne justifie
pas son ajout. Les analyses qualitative et d'erreur situent les
limites restantes : classes quasi vides et frontière avec le neutre.
